% titel: Blatt 0, A5: Ausmultiplizieren und Faktorisieren

% Teil I: Werkzeuge, kein Lösungsweg.

Terme ausmultiplizieren (Klammern auflösen) und Terme so weit wie möglich
faktorisieren.

\begin{werkzeugkasten}[Ausmultiplizieren]
Distributivgesetz: jeden Summanden der ersten Klammer mit jedem der zweiten
multiplizieren und addieren,
\begin{equation}
  (a_1+a_2+\dots)(b_1+b_2+\dots)=\sum_{i,j} a_i b_j .
\end{equation}
Nach dem Auflösen gleichnamige Terme zusammenfassen sowie die binomischen
Formeln nutzen: $(x+1)^2$, $(a-x)(a+x)=a^2-x^2$,
$(1-x)(1+x+x^2)$.
\end{werkzeugkasten}

\begin{werkzeugkasten}[Faktorisieren]
Umgekehrt Klammern bilden:
\begin{itemize}
  \item Zweite binomische Formel: $x^2+2x+1=(x+1)^2$.
  \item Dritte binomische Formel: $x^2-a^2=(x-a)(x+a)$.
  \item Nullstellen der quadratischen Form $ax^2+bx+c$: mit
        $x_{1,2}=\frac{-b\pm\sqrt{b^2-4ac}}{2a}$ ist die Faktorisierung
        $a(x-x_1)(x-x_2)$.
\end{itemize}
\end{werkzeugkasten}

\begin{achtung}[Darauf achten]
\begin{itemize}
  \item Vorzeichen beim Ausmultiplizieren von Differenzen (z. B.
        $y$ und $-y$ in $(x+1)(x-y+3)$).
  \item $x^2-a^2=(x-a)(x+a)$, aber $x^2+a^2$ ist über $\mathbb R$ nicht
        weiter zerlegbar.
  \item Beim Erkennen binomischer Formeln die Quadrate (z. B.
        $\tfrac14x^2=(\tfrac12x)^2$, $9=3^2$) und das doppelte Produkt
        prüfen.
\end{itemize}
\end{achtung}
